\chapter{Dataset Details and Preprocessing}
\label{AppendixB}
\lhead{Appendix B. \emph{Dataset Details and Preprocessing}}

\section{KAIST Multispectral Pedestrian Dataset}

The KAIST Multispectral Pedestrian Dataset~\cite{b_kaist} is a large-scale benchmark for multi-spectral pedestrian detection collected from an urban driving environment in South Korea. It contains over 95,000 colour-thermal image pairs captured under various illumination conditions (daytime, nighttime, dusk). The RGB and LWIR thermal cameras are calibrated and temporally synchronised.

\begin{itemize}
    \item \textbf{Original size:} $640 \times 480$ pixels
    \item \textbf{Preprocessed size:} $256 \times 256$ pixels (centre crop + resize)
    \item \textbf{Training pairs used:} 4,500 (stratified sample of 2,250 daytime + 2,250 nighttime)
    \item \textbf{Domain challenge:} High scene variability (>10 object categories), day/night illumination change, complex occlusion
\end{itemize}

\section{CBSR NIR Face Database}

The CBSR NIR Face Database provides near-infrared (NIR, 850 nm) and visible light face images. We treat NIR as proxy for thermal face imaging, as NIR-to-thermal mapping shares the cross-modal translation challenge (both are non-visible radiation domains).

\begin{itemize}
    \item \textbf{Modalities:} Near-infrared face $\to$ thermal face portrait
    \item \textbf{Preprocessed size:} $256 \times 256$ pixels (face-aligned crop)
    \item \textbf{Training pairs used:} 2,700
    \item \textbf{Domain challenge:} Subtle inter-subject thermal variation; consistent scene structure makes this the ``easiest'' domain
\end{itemize}

\section{Medical Knee Infrared Dataset}

The medical dataset consists of paired knee X-ray images and corresponding infrared thermography heatmaps indicating inflammation severity. This is a true cross-modal task: the input (X-ray) and target (pseudo-thermal heatmap) are fundamentally different imaging modalities.

\begin{itemize}
    \item \textbf{Modalities:} X-ray image $\to$ inflammation thermal heatmap
    \item \textbf{Preprocessed size:} $256 \times 256$ pixels
    \item \textbf{Training pairs used:} 4,500
    \item \textbf{Domain challenge:} X-ray and IR are physically unrelated modalities; the GAN must learn a highly non-trivial cross-modal mapping
\end{itemize}

\section{Agricultural Disease Datasets}

Two plant disease datasets from Kaggle were used, providing RGB leaf images with corresponding disease stress thermal maps.

\subsection{Tomato Leaf Disease}
\begin{itemize}
    \item \textbf{Modalities:} RGB tomato leaf $\to$ thermal vegetation stress map
    \item \textbf{Training pairs:} 663 (smallest dataset used)
    \item \textbf{Challenge:} Severe data scarcity; fine-grained disease hotspot boundaries
    \item \textbf{Disease classes:} Early blight, late blight, bacterial spot, healthy
\end{itemize}

\subsection{Chilli Pepper Disease}
\begin{itemize}
    \item \textbf{Modalities:} RGB bell/chilli pepper leaf $\to$ thermal disease footprint
    \item \textbf{Training pairs:} 2,227
    \item \textbf{Challenge:} Cross-class leaf shape variation; overlapping disease thermal signatures
    \item \textbf{Disease classes:} Bacterial spot, healthy
\end{itemize}

\section{Preprocessing Pipeline}

All datasets are preprocessed using the following standard pipeline:

\begin{enumerate}
    \item \textbf{Resize:} Bicubic interpolation to $256 \times 256$ pixels
    \item \textbf{Normalise RGB:} $[0, 255] \to [-1, 1]$ via $(x/127.5) - 1$
    \item \textbf{Normalise Thermal:} $[0, 255] \to [-1, 1]$ via $(y/127.5) - 1$
    \item \textbf{Thermal to Grayscale:} If multi-channel, convert using $y = 0.299R + 0.587G + 0.114B$ (luminance)
    \item \textbf{Train/Test Split:} 85\%/15\% random split, stratified where class labels available
    \item \textbf{Augmentation (training only):} Random horizontal flip ($p=0.5$)
\end{enumerate}

The preprocessing is implemented in \texttt{utils/dataset.py} using the \texttt{ThermalDataset} class, which inherits from \texttt{torch.utils.data.Dataset}.
